% main.tex - Plantilla de tesina (Universidad Veracruzana)
\documentclass[12pt,a4paper,oneside]{report}

% Codificación y español
\usepackage[utf8]{inputenc}
\usepackage[T1]{fontenc}
\usepackage[spanish]{babel}

% Márgenes según instrucción
\usepackage[a4paper,top=2.5cm,bottom=2.5cm,left=3cm,right=2.5cm]{geometry}

% Paquetes comunes
\usepackage{graphicx}
% Permite nombres de archivo de imágenes con varios puntos y espacios
\usepackage{grffile}
\usepackage{amsmath,amssymb}
\usepackage{booktabs}
\usepackage{caption}
\usepackage{subcaption}
\usepackage{float}
% Microtypografía: desactivar expansión si la distribución no tiene fuentes escalables
\usepackage[protrusion=true,expansion=false]{microtype}
\usepackage{hyperref}

\hypersetup{
  pdftitle={Título de la Tesina},
  pdfauthor={Yamileth Pérez Jaramillo},
  colorlinks=true, linkcolor=blue, citecolor=blue, urlcolor=blue
}

% Encabezados y pies
\usepackage{fancyhdr}
\pagestyle{fancy}
\fancyhf{}
\renewcommand{\headrulewidth}{0pt}
\fancyfoot[C]{\thepage}

% Metadatos rápidos
\newcommand{\titulo}{Título de la Tesina}
\newcommand{\autor}{Yamileth Pérez Jaramillo}
\newcommand{\facultad}{Facultad de Ingeniería Civil}

\begin{document}

% --- Portada y front matter ---
\pagenumbering{roman}
% portada.tex - Portada personalizada Universidad Veracruzana
\begin{titlepage}
\centering
\thispagestyle{empty}
\vspace*{1cm}
% Añade tu logo en figuras/logo_uv.png (se omite si no existe)
\IfFileExists{figuras/logo_uv.png}{%
	\includegraphics[width=3cm]{figuras/logo_uv.png}\par\vspace{1cm}
}{%
	\vspace{2.5cm}% espaciado alternativo si no hay logo
}
{\Large \facultad \par}
\vspace{1cm}
{\LARGE \textbf{UNIVERSIDAD VERACRUZANA}\par}
\vspace{1.5cm}
{\Huge\bfseries \titulo \par}
\vspace{1.5cm}
{\Large \autor\par}
\vfill
Trabajo de Tesina presentado para obtener el grado de\\[6pt]
Licenciatura en Ingeniería Civil\par
\vfill
{\large Asesor: Nombre del Asesor\par}
\vspace{1cm}
{\large Veracruz, México\\ 2026\par}
\end{titlepage}

\cleardoublepage

% dedicatoria.tex
\thispagestyle{empty}
\vspace*{6cm}
\begin{center}
\textit{Dedico esta tesina a...}
\end{center}

\cleardoublepage

% agradecimientos.tex
\chapter*{Agradecimientos}
\addcontentsline{toc}{chapter}{Agradecimientos}
Agradezco a todas las personas que brindaron apoyo técnico, académico y moral durante la realización de esta tesina.

\cleardoublepage

% resumen.tex
\chapter*{Resumen}
\addcontentsline{toc}{chapter}{Resumen}
\begin{abstract}
(En español) Resumen breve de la tesina: objetivos, metodología y resultados principales.
\end{abstract}

\chapter*{Abstract}
\addcontentsline{toc}{chapter}{Abstract}
\begin{abstract}
(In English) Short abstract: objectives, methodology and main results.
\end{abstract}

\cleardoublepage

\tableofcontents
\listoftables
\listoffigures
\cleardoublepage

% --- Contenido principal ---
\pagenumbering{arabic}
\setcounter{secnumdepth}{3}
\setcounter{tocdepth}{2}

% capitulo1.tex
\chapter{Introducción}
\section{Antecedentes}
El diseño de las placas base en estructuras de acero es un elemento crítico para la transferencia de cargas de la superestructura a la cimentación. Cambios en las disposiciones normativas y en los criterios de diseño (resistencia, ductilidad y verificación sísmica) pueden afectar dimensiones de placa, disposición y resistencia de pernos de anclaje, uso de rigidizadores y detalles constructivos.

En esta tesina se realiza una comparativa entre la NTC 2017 y la NTC 2023 aplicada al diseño de placas base para estructuras de acero, evaluando diferencias en hipótesis de carga, factores de reducción, prescripciones para anclajes y criterios de verificación. A continuación se muestran dos figuras que ayudan a contextualizar el problema y su origen sísmico.

\begin{figure}[H]
  \centering
  \includegraphics[width=0.4\textwidth]{\detokenize{figuras/001_Vista en corte de una conexión columna-cimentación columna de acero, placa base, pernos de anclaje, grout y pedestal de concreto..png}}
  \caption{Vista en corte de una conexión columna-cimentación columna de acero, placa base, pernos de anclaje, grout y pedestal de concreto.}
  \label{fig:baseplate}
\end{figure}

\begin{figure}[H]
  \centering
  \includegraphics[width=0.4\textwidth]{\detokenize{figuras/002_Mapa de la República Mexicana señalando el epicentro del sismo del 19-S 2017 y la Ciudad de México..png}}
  \caption{Mapa de la República Mexicana señalando el epicentro del sismo del 19-S 2017 y la Ciudad de México.}
  \label{fig:map2017}
\end{figure}

Las Figuras~\ref{fig:baseplate} y~\ref{fig:map2017} ilustran, respectivamente, el detalle constructivo objeto de estudio (placa base y pernos de anclaje) y el evento sísmico que motivó revisiones normativas posteriores que pueden influir en criterios de diseño.
\section{Objetivos}
\subsection{Objetivo general}
Enunciar el objetivo general.
\subsection{Objetivos específicos}
\begin{itemize}
  \item Objetivo específico 1.
  \item Objetivo específico 2.
\end{itemize}
\section{Justificación}
Justificación de la investigación y su aporte a la ingeniería civil.

% capitulo2.tex
\chapter{Marco Teórico}
\section{Conceptos fundamentales}
Desarrollo de definiciones y teorías relevantes.
\section{Normas y códigos aplicables}
Descripción de normas técnicas (ej.: ACI, NTC, AASHTO) y su aplicación.

% capitulo3.tex
\chapter{Estado del Arte}
\section{Revisión bibliográfica}
Resumen de trabajos previos y tendencias actuales.
\section{Comparación de métodos}
Tabla comparativa de técnicas y ventajas/desventajas.

% capitulo4.tex
\chapter{Metodología}
\section{Diseño de la investigación}
Descripción del enfoque metodológico aplicable a ingeniería civil.
\subsection{Trabajo de campo}
Muestreo, instrumentación y toma de datos en sitio.
\subsection{Procesos de laboratorio}
Ensayos normalizados y condiciones de ensayo.
\section{Modelado y análisis}
Modelado numérico, elementos finitos, parámetros y supuestos.
\subsection{Software y normas}
Listado sugerido: SAP2000, ETABS, Plaxis, OpenSees; normas de referencia.
\section{Validación y verificación}
Criterios de aceptación y procedimientos de verificación.

% capitulo5.tex
\chapter{Análisis y Resultados}
\section{Procesamiento de los datos}
Procedimientos estadísticos y filtrado de datos.
\section{Resultados}
Presentación de tablas, figuras y principales hallazgos.
\section{Discusión}
Interpretación técnica de los resultados y comparación con la literatura.


% --- Conclusiones y recomendaciones ---
% conclusiones.tex
\chapter*{Conclusiones y Recomendaciones}
\addcontentsline{toc}{chapter}{Conclusiones y Recomendaciones}
Aquí se plasman las conclusiones principales y recomendaciones para trabajos futuros.


% --- Anexos ---
\appendix
% anexo1.tex
\chapter{Anexo A: Instrumentos y Datos}
\section{Formularios y datos brutos}
Aquí se incluyen tablas, formatos y datos brutos que complementan la tesina.


% --- Bibliografía (BibTeX) ---
\bibliographystyle{unsrt}
\bibliography{bibliografia/references}

\end{document}
